\subsection{I-26. Genuine annihilators of the centered dense interpolation map}
\label{I-26}

\paragraph{Relation to prior work.}
Studying the vanishing ideal of the SV generator has an established
determinantal approach.  Hu, van Melkebeek, and Morgan (2024),
Theorem~1.3 (printed p.~6) and Corollary~1.5 (printed p.~7), give explicit
multilinear degree-$(k+1)$ annihilators of $\operatorname{SV}_k$
on $2k+1$ variables, settling the degree question discussed by
Fournier and Korwar (2018, WACT slides~13--19).
Their Theorems~1.9--1.10 retain a low-degree-monomial hypothesis
in the ROABP hitting/lower-bound assertion.  The construction
below follows this determinantal perspective, explicitly adds
invariance under the common translation, and proves an
invertible-factor ROABP representation.  No priority is claimed
for a vanishing-ideal construction in this setting.

\paragraph{Statement.}
Let $k\ge1$, let $n\ge3k+2$, and let a field $F$ contain
distinct Lagrange nodes $a_1,\ldots,a_n$.  For those nodes there
is an explicitly constructible nonzero homogeneous multilinear
polynomial $P\in F[X_1,\ldots,X_n]$ of degree $k+1$ such that
\[
 P\bigl(t\mathbf1+\operatorname{SV}_k(\mathbf y,\mathbf z)\bigr)=0.
 \tag{I.32}\label{eq:I-centered-sv-zero}
\]
It has a constant-boundary ROABP representation in every order
of width at most
\[
 W_k=\sum_{j=0}^{k+1}\binom{k+1}{j}^{\!2}
       =\binom{2k+2}{k+1},
 \tag{I.33}\label{eq:I-centered-sv-width}
\]
with commuting factors $I+X_iN_i$, where $N_i^2=0$.
Every interval of this representation is invertible, and its
constant matrix is the identity.  Thus, after constant
rank-one endpoints are added, $P$ is a genuine final rank-one
test whose interior satisfies all higher-rank preservation
requirements at zero.

In particular a centered $\operatorname{SV}_{k(w)}$ generator
valid for all widths at most $w$ and all $n$, over fields or
extensions supplying the required distinct nodes, must have
$W_{k(w)}>w$.  Equivalently, for each proposed $k$ the
counterexample requires an admissible $n\ge3k+2$.
Since $W_k\le4^{k+1}$, this implies
$k(w)=\Omega(\log w)$.  It does not rule out a logarithmic
number of pairs, and is not a lower bound for other generators.
All the general assertions hold in every characteristic under
the stated distinct-node hypothesis.

\paragraph{A determinant that removes the common center.}
Set $m=k+1$ and let
\[
 U_{i\ell}=a_\ell^i,\qquad 0\le i\le k,\quad1\le\ell\le n.
\]
In the $2k+2$ dimensional vector space of polynomials of degree
at most $2k+1$, impose the $k+1$ linear conditions
\[
 \sum_{\ell=1}^n a_\ell^i q(a_\ell)=0,
       \qquad0\le i\le k.
\]
Their common kernel has dimension at least $k+1$.
Choose linearly independent polynomials
$q_1,\ldots,q_m$ in that kernel by Gaussian elimination, and
put $V_{j\ell}=q_j(a_\ell)$.  Then $UV^T=0$.
Evaluation on the $n$ distinct nodes is injective on polynomials
of degree at most $2k+1$, since $n>2k+1$; hence $V$ has row
rank $m$.

Define
\[
 B(\mathbf X)=U\operatorname{diag}(X_1,\ldots,X_n)V^T,
 \qquad P(\mathbf X)=\det B(\mathbf X).
 \tag{I.34}
\]
Every maximal minor $U_S$ is a nonzero Vandermonde determinant.
At least one maximal minor $V_S$ is nonzero because $V$ has
full row rank.  Cauchy--Binet gives
\[
 P(\mathbf X)=\sum_{|S|=m}\det(U_S)\det(V_S)
                                  \prod_{\ell\in S}X_\ell.
\]
The nonzero minor of $V$ therefore supplies a nonzero coefficient.
This proves nonvanishing, homogeneity, and multilinearity in
every characteristic.

Because $UV^T=0$, one has $B(\mathbf X+t\mathbf1)=B(\mathbf X)$.
For every $i,j$, the polynomial $Y^i q_j(Y)$ has degree at most
$3k+1\le n-1$.  Exact Lagrange reproduction consequently gives
\[
 \sum_{\ell=1}^n a_\ell^i q_j(a_\ell)\lambda_\ell(Y)
    =Y^i q_j(Y).
\]
It follows that
\[
 B(\operatorname{SV}_k)
   =\sum_{s=1}^k z_s
       (1,y_s,\ldots,y_s^k)^T
       (q_1(y_s),\ldots,q_m(y_s)).
\]
This matrix is a sum of $k$ rank-one matrices but has size
$m=k+1$, so its determinant is identically zero.  The common
translation has no effect on $B$, proving (I.32).

\paragraph{An invertible, any-order ROABP of the stated width.}
For a generic $m\times m$ matrix $Z$, form the column vector
$\mathcal M(Z)$ of all square minors, including the empty
minor $1$.  There are $W_k$ such minors by the elementary
Vandermonde sum in (I.33).  They are linearly independent as
polynomials: distinct row/column index sets have distinct
row and column multidegrees, and each individual determinant
is nonzero.

For fixed vectors $u,v$ and an indeterminate $s$, expansion
of a minor under a rank-one update gives
\[
 \det(Z_{I,J}+s u_Iv_J^T)
 =\det Z_{I,J}
  +s\sum_{a\in I,\ b\in J}
        (-1)^{\operatorname{pos}_I(a)+\operatorname{pos}_J(b)}
        u_av_b\det Z_{I\setminus\{a\},J\setminus\{b\}}.
\]
Terms using two update columns vanish, since those columns are
proportional.  Thus there is a constant matrix $N(u,v)$ such that
\[
 \mathcal M(Z+suv^T)=(I+sN(u,v))\mathcal M(Z).
\]
Translation twice in the same direction and linear independence
of the minor polynomials imply $N(u,v)^2=0$ by comparing the
coefficient of two independent translation parameters.  Two
different rank-one translations commute, so their representing
matrices commute as well.  These arguments compare polynomial
identities and do not divide by factorials.

Let $u_\ell$ and $v_\ell$ be the $\ell$th columns of $U$ and
$V$, and set $N_\ell=N(u_\ell,v_\ell)$.  Starting at $Z=0$
and translating in these $n$ rank-one directions gives
\[
 \mathcal M(B(\mathbf X))
     =\prod_{\ell=1}^n(I+X_\ell N_\ell)\,\mathcal M(0).
\]
The coordinate of the full determinant is $P$.  The factors
commute and satisfy
$(I+X_\ell N_\ell)^{-1}=I-X_\ell N_\ell$.
This gives the advertised width-$W_k$ ROABP in any variable
order, with invertible factors and identity constants.
The starting row and final column may be taken as the
full-determinant selector and the empty-minor selector.
Adding their rank-one endpoint matrices produces an actual
rank-one test exactly as in I-21.  Every interior interval
has generic and constant rank $W_k$, so no missing rank
preservation hypothesis accounts for this annihilation.
\hfill$\square$

\paragraph{A width-four, one-pair example.}
Over $\mathbb Q$, use $n=5$ and nodes $0,1,2,3,4$.
The polynomials
\[
 q_1(Y)=Y^2-4Y+2,\qquad
 q_2(Y)=Y^3-\tfrac{77}{5}Y+\tfrac{54}{5}
\]
satisfy both kernel conditions.  Expanding (I.34) gives
$\det B=\frac65 Q$, where
\[
\begin{split}
 Q={}&3X_1X_2-4X_1X_3-15X_1X_4+16X_1X_5
       +4X_2X_3+8X_2X_4-15X_2X_5\\
    &+4X_3X_4-4X_3X_5+3X_4X_5.
\end{split}
\tag{I.35}
\]
Thus $Q(t\mathbf1+\operatorname{SV}_1)=0$ and $Q\ne0$.
For a smaller representation than the general width-six bound,
take row vectors and column vectors
\[
\begin{array}{c|ccccc}
i&1&2&3&4&5\\ \hline
a_i&(1,0)&(0,1)&(-28/97,-4/97)&(3/16,0)&(0,0)\\
b_i&(0,0)^T&(3,0)^T&(-4,4)^T&(-15,8)^T&(16,-15)^T.
\end{array}
\]
In blocks of sizes $1,2,1$, put
\[
 N_i=\begin{pmatrix}0&a_i&0\\0&0_{2\times2}&b_i\\0&0&0\end{pmatrix}.
\]
All triple products of these $N_i$ vanish, and
$e_1^TN_iN_je_4=a_i b_j$.
Consequently
\[
 e_1^T\prod_{i=1}^5(I_4+X_iN_i)e_4
      =\sum_{i<j}(a_ib_j)X_iX_j=Q.
\]
The factors are upper triangular with diagonal entries one,
so every interior interval remains invertible at every point.
This refutes the one-pair proposal already at width four;
it does not assert that width four is the smallest possible
width of a centered one-pair annihilator of arbitrary degree.

\paragraph{A concrete repair for this family, and its scope.}
The same counterexamples are detected by a different simple
map.  Since every $U_S$ is nonzero, the support of $P$ is
exactly the set of bases of the column matroid of $V$.
Pairwise distinct element weights have a unique minimum-weight
matroid basis, by the greedy exchange property.  Hence
\[
 P(T,T^2,\ldots,T^n)\ne0,
\]
with degree at most $n(k+1)$.
This is an instance of the basis argument in I-03 and gives a
polynomial scalar sampling grid for this particular family:
choose $n(k+1)+1$ distinct scalar values, in an extension if
the ground field is too small.  This additional grid-size
condition is separate from the $n$ distinct Lagrange nodes.
It demonstrates that the dense-SV counterexample is a precise
test of that map, not a general obstruction to cocharacter
sampling.  The open dense-map question is whether a
logarithmic or polynomial number of pairs suffices for the
full rank-test class, or whether a different structured family
can handle its cancellations with polynomial dependence on
width.  Neither assertion is proved by this example.
